\chapter{An Executable Modern Training Core}
\label{ch:modern-core}

The earlier chapters teach TinyGPT by taking its components apart. This chapter
puts a complete, deliberately small model under test. A reader should be able to
change an equation, see a failing test, and explain why the test failed. That is
the bridge from understanding a model to engineering a training system.

The implementation lives in \texttt{code/minillm/}. It uses PyTorch autograd and
CPU FP32 arithmetic; it does not implement a new differentiation framework.
Its name, \term{TinyLM3}, identifies a modern teaching lineage, not a compatible
implementation of a commercial model. Nothing downloads weights or a dataset.
The reference run needs neither a GPU nor an external service.

\section{Four Planes, One Correctness Contract}

\begin{figure}[htbp]\centering
\begin{tikzpicture}
\node[engineop,text width=11.8cm] (model) at (0,0) {01 / MODEL SEMANTICS\quad Tokens, causal logits, objective};
\node[enginevalue,text width=11.8cm] (execute) at (0,-1.65) {02 / EXECUTION\quad Tensors, autograd, optimizer, precision, memory};
\node[engineop,text width=11.8cm] (distribute) at (0,-3.3) {03 / DISTRIBUTION\quad Ownership, collectives, shards, pipeline schedule};
\node[enginestate,text width=11.8cm] (operate) at (0,-4.95) {04 / OPERATIONS\quad Data identity, checkpoint, recovery, provenance};
\draw[enginearrow] (model) -- node[right,enginenote]{realize} (execute);
\draw[enginearrow] (execute) -- node[right,enginenote]{partition} (distribute);
\draw[enginearrow] (distribute) -- node[right,enginenote]{recover} (operate);
\node[enginenote] at (0,-6) {Correctness crosses every plane. This milestone executes planes 1, 2 and local recovery.};
\end{tikzpicture}

\caption{A training system has four distinct responsibilities. Distribution is
a planned extension here, not an implemented feature of the CPU trainer.}
\label{fig:training-planes}
\end{figure}

The model specifies which function is differentiated. Execution decides how
tensors, gradients and optimizer state realize that function. Distribution
partitions that work across processes and devices. Operations makes the run
identifiable, recoverable and observable. A correct local model does not prove
a correct collective; a successful checkpoint write does not prove a correct
restart. We will add each claim only with its own test.

\section{An Explicit Lineage Change}

Do not silently replace the definitions in the introductory chapters. Their
learned position embeddings, LayerNorm and GELU remain a useful TinyGPT path.
The modern core makes a different, explicit choice:

\begin{center}
\begin{tabular}{lll}\toprule
Decision & Earlier TinyGPT teaching path & TinyLM3 executable core \\\midrule
Position & Learned embedding table & Adjacent-pair RoPE on Q and K \\
Normalization & LayerNorm & Pre-block RMSNorm and final RMSNorm \\
Attention & Multi-head & Grouped-query, causal \\
MLP & GELU two-layer MLP & SwiGLU gate, up and down \\
Output weights & Tying discussed as an option & Untied embedding and head \\
Dropout & Explained and configurable in examples & Absent in this milestone \\\bottomrule
\end{tabular}
\end{center}

The new choices follow RMSNorm, rotary position embeddings, GQA and gated MLP
literature \cite{zhang2019rmsnorm,su2021rope,ainslie2023gqa,shazeer2020glu}.
They do not imply a quality improvement on our eight-character pattern. This
fixture tests mechanics, not language competence.

\section{The Typed Forward Pass}

Let \(B\) denote sequences per batch, \(T\) positions per sequence, \(V\)
vocabulary size, \(D\) residual width, \(H_q\) query heads, \(H_{kv}\) key/value
heads and \(d=D/H_q\) head width. Configuration rejects nonpositive dimensions,
nondivisible head groups, odd rotary widths and invalid epsilon. The default
is \(V=8,T\leq16,D=32,L=2,H_q=4,H_{kv}=2,F=64\), where \(F\) is MLP width.

The parameter count is \term{DERIVED}, with untied vocabulary matrices and no biases:
\[
 P=2VD+L\bigl(2D^2+2DH_{kv}d+3DF+2D\bigr)+D=19{,}104.
\]
The terms count embedding/head, Q/O, K/V, three MLP matrices, two block
normalization gains and final gain. A test checks this formula against every
registered parameter; a different weight-tying policy would change the count.

The data layer consumes \(T+1\) consecutive int64 token IDs:
\[
 x_{b,t}=s_{b,t},\qquad y_{b,t}=s_{b,t+1},\qquad 0\leq t<T.
\]
Targets are already shifted. Shifting them again changes the task into
two-token-ahead prediction. Taking \(y=x\) makes a copying shortcut possible.
The final-window test uses exactly \(T+1\) tokens, a boundary that catches an
off-by-one error in the original batch example.

Embedding lookup maps \([B,T]\) IDs to \([B,T,D]\) activations. The embedding
table is a persistent parameter, while the batch activation belongs to this
forward pass. Each block preserves residual width; the final head maps to raw
logits \(Z\in\mathbb{R}^{B\times T\times V}\). Do not apply softmax before
\texttt{cross\_entropy}: the library combines normalization and logarithms
using a numerically stable formulation.

\begin{figure}[htbp]\centering
\begin{tikzpicture}
\node[enginevalue] (input) at (0,0) {Residual input\\\(X:[B,T,D]\), FP32};
\node[engineop] (norm1) at (0,-1.7) {RMSNorm\\same shape};
\node[engineop] (attn) at (0,-3.4) {Q/K RoPE + causal GQA\\\(H_q=4,\ H_{kv}=2\)};
\node[enginevalue] (sum1) at (0,-5.1) {Add residual\\\(X'=X+A\)};
\node[engineop] (norm2) at (5,-5.1) {RMSNorm + SwiGLU\\hidden width \(F\)};
\node[enginevalue] (sum2) at (5,-6.8) {Add residual\\\(X''=X'+M\)};
\draw[enginearrow] (input) -- (norm1);
\draw[enginearrow] (norm1) -- (attn);
\draw[enginearrow] (attn) -- (sum1);
\draw[enginearrow] (sum1) -- (norm2);
\draw[enginearrow] (norm2) -- (sum2);
\draw[enginearrow] (input.west) -- ++(-1.2,0) |- node[pos=.25,left,enginenote]{identity} (sum1.west);
\draw[enginearrow] (sum1.south) |- node[pos=.7,below,enginenote]{identity} (sum2.west);
\node[enginenote,text width=4cm] at (5,-1.3) {Persistent parameters:\\Q, K, V, O; gate, up, down;\\two learned norm gains};
\node[enginenote,text width=4cm] at (5,-3.2) {Autograd retains needed\\intermediates until backward.\\No inference KV cache.};
\end{tikzpicture}

\caption{One pre-normalized TinyLM3 block. Solid arrows carry activations;
the two bypasses are additive identity paths, not additional learned projections.}
\label{fig:modern-block}
\end{figure}

\section{RMSNorm: A Small Formula With a Real Gradient}

For one token vector \(x\in\mathbb{R}^{D}\), learned gain
\(g\in\mathbb{R}^{D}\) and positive scalar \(\epsilon\), define
\[
 r=\left(\frac{1}{D}\sum_{j=1}^{D}x_j^2+\epsilon\right)^{-1/2},
 \qquad u_i=g_i x_i r.
\]
There is no subtraction of the mean. Gain is applied after the shared scale,
and the output need not have unit RMS. For \(x=(1,-2,3,-4)\), the mean square
is \(7.5\); the test compares every coordinate against
\(x/\sqrt{7.5+10^{-5}}\) with unit gain.

To see why backward couples coordinates, differentiate the shared scale:
\[
 \frac{\partial u_i}{\partial x_k}
 =g_i\left(r\,\delta_{ik}-\frac{x_i x_k}{D}r^3\right).
\]
The second term connects every input coordinate to every output coordinate.
\texttt{gradcheck} perturbs the double-precision input and compares finite
differences with autograd. This is a test of the implemented local derivative,
not proof against every floating-point overflow case. Mixed precision and
stabilized large-magnitude reductions belong to the next milestone.

\section{RoPE and Head Grouping}

For adjacent coordinates \((a,b)\) at position \(p\) and frequency
\(\omega_j=10000^{-2j/d}\), RoPE applies
\[
 \begin{bmatrix}a'\\b'\end{bmatrix}
 =\begin{bmatrix}\cos(p\omega_j)&-\sin(p\omega_j)\\
 \sin(p\omega_j)&\cos(p\omega_j)\end{bmatrix}
 \begin{bmatrix}a\\b\end{bmatrix}.
\]
The transform preserves a pair's squared length in exact arithmetic.
Position zero is the identity. Tests check both properties and the gradient.
The implementation rotates Q and K, never V. It resets positions to zero for
each training window; incremental decoding with position offsets is not present.

Q has shape \([B,H_q,T,d]\), while K and V have shape
\([B,H_{kv},T,d]\) before head expansion. With group size
\(G=H_q/H_{kv}\), query head \(h\) reads KV head \(\lfloor h/G\rfloor\).
Our readable implementation physically repeats K/V along the head dimension.
This makes grouping easy to inspect but is not a memory-optimal kernel.

For one query head, the mathematical attention result is
\[
 A=\operatorname{softmax}\!\left(\frac{QK^{\mathsf T}}{\sqrt d}+M\right)V,
 \qquad
 M_{t,s}=\begin{cases}0&s\leq t,\\-\infty&s>t.\end{cases}
\]
Every row retains its diagonal entry, so no row is fully masked. The score and
probability matrices explicitly occupy \([B,H_q,T,T]\); this is not
FlashAttention. An independent test loops over each head and visible position,
computes dot products and weighted sums, then compares with the vectorized
implementation. Another changes only future token IDs and demands unchanged
prefix logits. Neither test can be replaced by a falling loss curve.

\section{SwiGLU and the Objective}

With row-oriented activations and matrices
\(W_g,W_u\in\mathbb{R}^{D\times F}\),
\(W_d\in\mathbb{R}^{F\times D}\), the MLP is
\[
 \operatorname{MLP}(X)=
 \bigl(\operatorname{SiLU}(XW_g)\odot XW_u\bigr)W_d,
 \qquad \operatorname{SiLU}(z)=z\,\sigma(z).
\]
PyTorch stores each linear weight as output-by-input, so its physical shapes
are transposed relative to this equation. Checkpoint conversion must preserve
that distinction. All three projections are bias-free.

Mean next-token cross entropy is
\[
 \mathcal L=-\frac{1}{BT}\sum_{b=1}^{B}\sum_{t=1}^{T}
 \log\frac{\exp Z_{b,t,y_{b,t}}}{\sum_{v=1}^{V}\exp Z_{b,t,v}}.
\]
One whole-model finite-difference test perturbs an output-head parameter and
checks its loss derivative. It complements operator tests by exercising the
complete route from token IDs to the scalar objective.

\section{One Optimizer Boundary}

\begin{figure}[htbp]\centering
\begin{tikzpicture}
\node[enginevalue] (data) at (0,0) {Data batch\\\(x_t,y_t:[B,T]\), int64};
\node[engineop] (forward) at (4.5,0) {Forward + cross entropy\\scalar mean loss};
\node[engineop] (back) at (9,0) {Backward\\parameter gradients};
\node[enginestate] (param) at (0,-2.5) {Persistent model\\\(\theta_t\)};
\node[enginestate] (opt) at (4.5,-2.5) {Persistent AdamW state\\\(m_t,v_t,\text{step}_t\)};
\node[engineop] (update) at (9,-2.5) {Clip, step, clear gradients\\\(\theta_{t+1}\)};
\draw[enginearrow] (data) -- (forward);
\draw[enginearrow] (forward) -- (back);
\draw[enginearrow] (param.north east) -- node[above,enginenote]{read weights} (forward.south west);
\draw[enginearrow] (back) -- (update);
\draw[enginearrow] (opt) -- (update);
\draw[enginearrow] (update.south) -- ++(0,-.7) -| node[pos=.25,below,enginenote]{new state for the next optimizer boundary} (param.south);
\node[enginenote,text width=11cm] at (4.5,-5.1) {Saved activations: forward to backward. Gradients: backward to clear.\\Parameters and moments: across steps. Checkpoint: only after update completes.};
\end{tikzpicture}

\caption{State lifetime is part of correctness. Saved activations, gradients
and optimizer moments have different owners and release boundaries.}
\label{fig:step-lifetime}
\end{figure}

The trainer clears gradients, computes token-weighted microbatch losses,
backpropagates, rejects nonfinite gradients, clips the total gradient norm to
one, performs AdamW, and clears gradients again. AdamW decouples weight decay
from the moment estimates \cite{loshchilov2019adamw}. The pinned reference uses
\texttt{foreach=False} and decays every parameter, including norm gains. This
is an explicit teaching policy, not a recommended production hyperparameter recipe.

If microbatch \(i\) contains \(n_i\) targets, accumulate
\[
 \mathcal L=\sum_i \frac{n_i}{\sum_j n_j}\mathcal L_i.
\]
Dividing each mean loss by the number of microbatches is equivalent only when
their target counts match. The test deliberately splits a batch into sizes one
and two. Compare parameter updates within FP32 tolerance, since changing the
reduction order can change final bits. The optimizer test separately checks
the first AdamW update against its scalar formula.

\section{Two Sanity Experiments, Different Questions}

First, fit a deterministic repeating eight-character cycle. A large decrease
in loss shows that the model, backward pass and optimizer are connected. It
does not establish language generalization: train and evaluation windows share
the same simple rule, and the model can memorize it.

Second, generate independent uniform token streams. Conditional on any prefix,
the next token still has probability \(1/V\), so its irreducible expected
cross entropy is \term{DERIVED}:
\[
 H(Y\mid X)=\log V=\log 8\approx2.07944\ \text{nats per token}.
\]
Use separate seeded generators for fresh training tokens and held-out evaluation
tokens. Repeatedly fitting one finite random dataset instead tests memorization.
The test allows a 0.12-nat interval around entropy after 100 updates and evaluates
4096 independent target positions. This deterministic regression threshold is
not a universal statistical confidence interval. A much lower held-out loss
would be suspicious: inspect target alignment, causal masking and split leakage.

\labtitle{Run, Break, Explain}
\begin{lstlisting}[language=bash]
python3 -m venv .venv
.venv/bin/python -m pip install -r requirements.txt
make check PYTHON=.venv/bin/python
make smoke PYTHON=.venv/bin/python
\end{lstlisting}

Remove the causal mask and rerun the future-token test. Change the batch upper
bound and rerun the final-window test. Replace token-weighted accumulation with
equal microbatch weighting and observe its regression. Restore each change
before proceeding. Record the actual command, dependency version, seed, device
and measured loss; do not copy someone else's output and call it a measurement.

\checkpoint
\begin{enumerate}
\item Why can a model memorize random training data but not predict fresh random targets?
\item Which RMSNorm derivative term couples different coordinates?
\item Why does repeating KV heads make this reference inefficient but still mathematically useful?
\item Which test fails if targets are shifted twice? Design an additional direct oracle.
\item What state survives an optimizer step, and what state may be released?
\end{enumerate}
